\section{Conclusiones Parciales}

Los resultados preliminares obtenidos hasta el momento permiten establecer las siguientes conclusiones parciales:

La detección binaria de archivos cifrados por ransomware mediante métricas estadísticas es viable, alcanzando hasta 100\% de accuracy con modelos basados en árboles de decisión. Esto confirma la efectividad de métricas como la entropía de Shannon, Chi-Cuadrado, estimación Monte Carlo y el coeficiente de correlación serial de bytes para la tarea de identificar si un archivo ha sido encriptado. No obstante, como señalan Lee et al.~\cite{lee2022}, la entropía de Shannon por sí sola no es la métrica óptima y debe combinarse con métricas complementarias para reducir falsos positivos.

La clasificación multiclase de familias de ransomware basada exclusivamente en propiedades estadísticas de archivos cifrados no es viable, alcanzando un máximo del 15,4\% de accuracy sobre 31 clases. Este resultado se explica por la convergencia estadística inherente a los algoritmos criptográficos modernos, que producen salidas indistinguibles entre familias~\cite{paper_4_comparison_of_entropy}.

La evaluación de herramientas existentes revela que las notas de rescate son más discriminativas que los archivos cifrados: ID Ransomware~\cite{idransomware} identifica el 71,93\% de las familias con notas versus el 66,67\% con archivos. Esto motiva el siguiente paso de la investigación.

\section{Trabajo Pendiente}

El trabajo pendiente comprende:

\begin{enumerate}
    \item \textbf{Construcción del corpus de notas:} Recopilar notas de rescate para las 30 familias de NapierOne desde los repositorios ThreatLabz~\cite{threatlabz}, Lemmou~\cite{paper_5_dataset} y fuentes de seguridad públicas, asegurando un mínimo de 2 notas por familia.

    \item \textbf{Clasificación NLP:} Implementar y evaluar el clasificador basado en TF-IDF con n-gramas de palabras y caracteres, comparando LinearSVC, Regresión Logística, Random Forest y KNN.

    \item \textbf{Ejecución en NapierOne-Small:} Repetir los experimentos de clasificación binaria y multiclase sobre el dataset NapierOne-Small (1.000 archivos por familia) para obtener resultados más robustos con las 6 métricas estadísticas completas.

    \item \textbf{Análisis comparativo:} Comparar cuantitativamente la efectividad de ambos enfoques (archivos cifrados vs. notas de rescate) para la identificación de familias.

    \item \textbf{Redacción final:} Completar el capítulo de resultados con los datos experimentales definitivos y las conclusiones finales del trabajo.
\end{enumerate}
